\subsection{AC-G82: effective fresh quadratic hosts under A}
\label{sec:cubic-effective-fresh-host}
\paragraph{Scope and notation.}
Here \(p=269\), \(K=\mathbb Q(\sqrt{-3})\), and A has exactly its stipulated meaning: the Riemann hypothesis for \(\zeta\) and every primitive Dirichlet \(L\)-function. Let \(S\) be a finite set of rational primes containing \(2,3,p\), and put
\[
 M=8\prod_{\substack{\ell\in S\\\ell\ {\rm odd}}}\ell,\qquad
 B(x)=2(x\log x)^2\quad(x\geq3).
\]
If an old Galois extension \(\Omega/\mathbb Q\) is being retained, assume that it contains \(K\) and is unramified outside \(S\). It may be infinite. The statements below quantify the fresh quadratic choices in AC-G77 and AC-G81. They do not supply the different nonabelian Frobenius approximation used to choose the level-raising prime.

\paragraph{Published input, with its exact scope.}
Bach--Sorenson, \emph{Explicit bounds for primes in residue classes}, Math.\ Comp.\ 65 (1996), 1717--1735, Theorem 5.3, printed pp.\ 1731--1732, gives, under its ERH assumption, a prime in any reduced residue class modulo \(m\) below \(B(m)\). For this application the field is \(\mathbb Q(\mu_m)\); the assumption at the end of their \S2, printed p.\ 1720, is zero-freeness of its zeta function and that of \(\mathbb Q\) in \(\Re s>1/2\). The factorization of the cyclotomic zeta function into primitive Dirichlet \(L\)-functions shows that A supplies precisely this input. We retain the published theorem, including its numerical estimates, as background. Its whole printed Theorem 5.3 proof was read; its numerical program was not rerun.

\paragraph{Finding AC-G82.}
There are distinct primes \(r_0,r_1\notin S\) such that
\[
 r_0\equiv-1\pmod M,\qquad r_0<B(M),\qquad
 r_1\equiv-1\pmod{Mr_0},\qquad r_1<B(Mr_0).
 \tag{G82.1}
\]
Write
\[
 D_i=\mathbb Q(\sqrt{-r_i}),\quad
 F=KD_1,\quad F^+=\mathbb Q(\sqrt{3r_1}),\quad
 L=KD_0D_1,\quad L^+=\mathbb Q(\sqrt{3r_0},\sqrt{3r_1}).
\]
The quadratic character of \(D_0/\mathbb Q\) is odd, is ramified only at \(r_0\), and is trivial on every decomposition group at a prime of \(S\). The same local triviality holds for \(D_1\) at every prime of \(S\cup\{r_0\}\).
The fields \(F,L\) are CM of degrees \(4,8\); the displayed fields \(F^+,L^+\) are their maximal real subfields. Both CM extensions are unramified at all finite places. Their absolute discriminants are
\[
 |\operatorname{Disc}F|=9r_1^2,\quad
 \operatorname{Disc}F^+=3r_1,\quad
 |\operatorname{Disc}L|=81r_0^4r_1^4,\quad
 \operatorname{Disc}L^+=9r_0^2r_1^2.
 \tag{G82.2}
\]
Every prime of \(F^+\) over \(S\cup\{r_0\}\) splits in \(F\), and every prime of \(L^+\) over \(S\) splits in \(L\). At primes over \(S\), each local extension \(L/F\) has degree one. In particular the added quadratic base change preserves the local field at \(269\).
Moreover,
\[
 F\cap\Omega=K,\qquad L\cap\Omega=K.
 \tag{G82.3}
\]
There is the explicit estimate
\[
 |\operatorname{Disc}L|
   <81\,2^{32}(M\log M)^{32}.
 \tag{G82.4}
\]
For a single fresh G77 quadratic, choosing \(r\equiv-1\pmod M\) gives instead
\[
 |\operatorname{Disc}(K\mathbb Q(\sqrt{-r}))|
   <36(M\log M)^4.
 \tag{G82.5}
\]

\paragraph{Proof.}
Apply the published least-prime theorem first to the reduced residue class \(-1\) modulo \(M\), and then to \(-1\) modulo \(Mr_0\). Since every old odd prime divides \(M\), neither choice can be an old prime; both are odd, and the second cannot equal \(r_0\). This proves (G82.1). Both are \(7\) modulo \(8\) and \(2\) modulo \(3\).
The negative squarefree integer \(-r_i\) is \(1\) modulo \(4\), so \(D_i\) has fundamental discriminant \(-r_i\), and its quadratic character is odd. For an odd \(\ell\mid M\), the unit \(-r_i\) is \(1\) modulo \(\ell\), hence is a square in \(\mathbb Q_\ell\) by Hensel's lemma. At \(2\), it is \(1\) modulo \(8\), which is the square criterion for odd \(2\)-adic units. The same reasoning at \(r_0\) applies to \(D_1\). These prove every asserted local triviality, not just unramifiedness.

The square classes \(-3,-r_0,-r_1\) are independent over \(\mathbb Q\), as their valuations at \(3,r_0,r_1\) show. Thus the degrees are as stated. Complex conjugation changes the sign of all three chosen square roots. Its fixed fields are exactly the displayed real fields, proving the CM and maximal-real assertions.
For \(F\), the three quadratic subfields have fundamental discriminants \(-3,-r_1,3r_1\). The conductor--discriminant formula gives \(9r_1^2\). For \(L^+\), the nontrivial quadratic subfields have fundamental discriminants \(3r_0,3r_1,r_0r_1\), whose product is \(9r_0^2r_1^2\). For \(L\), the seven discriminants are
\[
 -3,\ -r_0,\ -r_1,\ 3r_0,\ 3r_1,\ r_0r_1,\ -3r_0r_1.
\]
Each is squarefree and \(1\) modulo \(4\), so their absolute product is \(81r_0^4r_1^4\). These are instances of the conductor--discriminant identity, for example Bach--Sorenson \S2, equation (2.3), printed p.\ 1720. The discriminant tower formula now gives relative discriminant of norm one for both \(F/F^+\) and \(L/L^+\), since their absolute discriminants are respectively the squares of the real-field discriminants. Thus both are unramified at every finite place.

A quadratic field \(D_1\) splitting at \(\ell\) gives \(F=F^+D_1\) split above every place over \(\ell\); here \(F=F^+D_1\) follows directly from the square roots. This proves the claimed splitting for \(F/F^+\). Likewise \(L=L^+D_1\), proving it for \(L/L^+\). Finally \(L=FD_0\), and \(D_0\) splits over \(\mathbb Q_\ell\) at each \(\ell\in S\), so all the local degrees in \(L/F\) there are one.

For the intersection assertions, \(D_1\) is ramified at the fresh prime \(r_1\), whereas \(\Omega\) is unramified there. Hence \(D_1\not\subset\Omega\), and \(F\cap\Omega=K\) because \([F:K]=2\) and \(K\subset\Omega\). The group \(\operatorname{Gal}(L/K)\) is \(C_2^2\). Its three quadratic intermediate fields over \(K\) are \(KD_0,KD_1,K\mathbb Q(\sqrt{r_0r_1})\). Each is ramified over \(K\) at \(r_0\) or \(r_1\): \(K/\mathbb Q\) is unramified at those odd primes, so adjoining a square root with odd valuation is ramified. None can be contained in \(\Omega\). A nontrivial subextension of a \(C_2^2\)-extension contains a quadratic subextension, proving \(L\cap\Omega=K\). This argument applies unchanged to an infinite \(\Omega\), since the intersection is finite.

Set \(t=\log M\). Since \(M\geq8\cdot3\cdot269=6456\), we have \(t>8\), and
\[
 \log(2M^3t^2)=\log2+3t+2\log t\leq4t.
\]
For completeness, \(t-2\log t-\log2\) is increasing for \(t>2\) and positive at \(t=8\). From (G82.1),
\[
 r_0<2M^2t^2,\qquad
 r_1<2(Mr_0\log(Mr_0))^2
       <128M^6t^6.
\]
Consequently \(r_0r_1<256M^8t^8\), and (G82.2) gives (G82.4). A single choice \(r<B(M)\) gives (G82.5) by \(9r^2<36M^4t^4\). This completes the proof.

\paragraph{What this supplies and what remains.}
The construction is effective: test successive integers in the two displayed progressions for primality; the proved bounds guarantee termination. No such enumeration is claimed to have been executed during the outage. In AC-G81 it makes the fresh scalar-witness prime locally trivial at every old bad place and bounds both fresh quadratic discriminants. It does not bound the auxiliary prime obtained by approximating a nonabelian Hida-family Frobenius condition. It also does not control newly acquired class sectors or verify the numerical \(\mathrm{Frame.Ok}\), \(\mathrm{EC}\), \(\mathrm{Big}\), counting, or per-curve \(J\)-containment conditions after host enlargement.

\paragraph{Geometric dictionary.}
The fresh quadratic covers are finite etale over \(\operatorname{Spec}\mathbb Z[1/r_i]\), with specified split fibers over the old prime set. Their normalized CM covers of the maximal real rings are finite etale at every finite place, by the relative-discriminant computation. The intersection calculation says that these degree-two covers do not become constant in the old extension. This uses the already cited finite-etale/Galois correspondence, SGA1, Expose V, \S\S4 and 7; no new archival cote is asserted.
